% A heap block: one header word, then the fields. The pointer lands
% AFTER the header. Used in M10-L02.
\documentclass[tikz,border=4pt]{standalone}
\usepackage{tikz}
\input{_m10-styles.texinput}
\begin{document}
\begin{tikzpicture}[m10, x=1mm, y=-1mm, font=\large]
  % header word
  \node[stackcell, minimum width=54mm, minimum height=14mm, font=\large] (hdr) at (27,0)
    {header};
  \node[note, anchor=north, font=\large] at (27,9) {size \(\cdot\) tag \(\cdot\) GC colour};

  % field words
  \node[heapcell, minimum width=26mm, minimum height=14mm, font=\large] (f0) at (67,0) {field 0};
  \node[heapcell, minimum width=26mm, minimum height=14mm, font=\large] (f1) at (94,0) {field 1};
  \node[heapcell, minimum width=26mm, minimum height=14mm, font=\large] (f2) at (121,0) {field 2};

  % the pointer points just past the header
  \draw[flow, line width=1pt] (67,-20) -- (f0.north west);
  \node[note, anchor=south, font=\large] at (67,-21) {the pointer};

  \node[note, anchor=north west, font=\large\itshape] at (0,16)
    {The header is one hidden word: it is why \texttt{Array.length} is \(O(1)\).};
\end{tikzpicture}
\end{document}
